\chapter{Karp--Rabin}
\label{cap:karp-rabin}

\section{Coda: le applicazioni del suffix tree}

\scan{26}
Un'ultima annotazione chiude il blocco sui suffix tree. Il suffix tree è la struttura che
realizza il \emph{pre-processing del testo}; Gusfield vi dedica due capitoli di applicazioni,
che si dividono in due famiglie:
\begin{itemize}
  \item applicazioni \emph{semplici};
  \item applicazioni \emph{non semplici}, che nascono dalla combinazione
        ``suffix tree $+$ Harel--Tarjan''.
\end{itemize}

Il motivo della seconda voce è che, una volta costruito il suffix tree di $T$, moltissime
domande su $T$ si riducono a calcolare il $\lca$ di due foglie. Il teorema di Harel--Tarjan
garantisce che, dopo un pre-processing lineare nella dimensione dell'albero, ogni query di
lowest common ancestor viene risolta in tempo $\Oh(1)$: è questo che rende praticabili --- e
al tempo stesso non banali --- le applicazioni della seconda famiglia. Al teorema, e alla sua
dimostrazione, è dedicato un capitolo a sé più avanti.

\begin{figure}[H]\centering% stree-lca-montagna -- pag. 26
% Il suffix tree disegnato "a montagna": radice in cima, foglie sulla linea di base,
% il nodo lca(l1,l2) a meta' altezza raggiunto in O(1) grazie a Harel--Tarjan.
\begin{tikzpicture}[
    >=Latex,
    ramo/.style={draw=black!75,line width=0.5pt},
    nodo/.style={circle,fill=black!80,inner sep=0pt},
    et/.style={font=\footnotesize}
  ]

  % ---- intestazione dello schizzo -------------------------------------------
  \node[font=\small] at (0.1,4.15) {\STree $+$ \underline{Harel--Tarjan}};

  % ---- sagoma "a montagna" del suffix tree ----------------------------------
  \draw[line width=0.9pt,draw=black!55,fill=black!5]
      (0.40,3.50) -- (0.10,2.85) -- (-0.15,2.60) -- (-0.45,2.35)
      -- (-0.85,1.95) -- (-1.05,1.45) -- (-1.55,1.05) -- (-1.85,0.60)
      -- (-2.45,0.02) -- (2.45,0.02)
      -- (2.00,0.55) -- (1.80,1.05) -- (1.35,1.45) -- (1.20,2.05)
      -- (0.85,2.35) -- (0.70,2.80) -- cycle;

  % ---- "terreno": e' il livello delle foglie --------------------------------
  \draw[line width=0.9pt,draw=black!55]
      plot[smooth,tension=0.8] coordinates
      {(-2.90,0.12) (-2.30,-0.02) (-1.60,0.07) (-0.90,-0.02)
       (-0.10,0.07) (0.70,-0.02) (1.50,0.07) (2.20,-0.02) (2.90,0.12)};

  % ---- cammino radice -> lca (nel quaderno e' ripassato due volte) ----------
  \draw[line width=2.4pt,black!20] (0.40,3.50) -- (0.10,1.85);
  \draw[ramo,dashed] (0.40,3.50) -- (0.10,1.85);

  % ---- archi lca -> foglie --------------------------------------------------
  \draw[ramo] (0.10,1.85) -- (-0.55,1.05) -- (-0.95,0.02);
  \draw[ramo] (0.10,1.85) -- (0.35,0.02);

  % ---- nodi -----------------------------------------------------------------
  \node[nodo,minimum size=3pt] (rad) at (0.40,3.50) {};
  \node[nodo,minimum size=6pt] (a)   at (0.10,1.85) {};
  \node[nodo,minimum size=4pt] (b)   at (-0.55,1.05) {};
  \node[nodo,minimum size=4pt] (l1)  at (-0.95,0.02) {};
  \node[nodo,minimum size=4pt] (l2)  at (0.35,0.02) {};

  % ---- etichette ------------------------------------------------------------
  \node[et,right=3pt] at (rad.east) {radice};
  \node[font=\small\itshape] at (-1.30,0.62) {\STree};

  \draw[black!45,line width=0.4pt] (a) -- (2.10,2.45);
  \node[et,anchor=west] at (2.10,2.45) {$\lca(\ell_1,\ell_2)$};

  \node[et,below=2pt] at (l1) {$\ell_1$};
  \node[et,below=2pt] at (l2) {$\ell_2$};

\end{tikzpicture}

\caption{Applicazioni ``non semplici'': dentro il suffix tree di $T$ si individua il lowest
common ancestor di due foglie, che con il pre-processing di Harel--Tarjan costa $\Oh(1)$.}
\label{fig:stree-lca-montagna}
\end{figure}

\section{Impostazione: dalle stringhe ai numeri}

Siano $T,P \in \S^{*}$ e poniamo $m = \abs{P}$, $n = \abs{T}$. Negli appunti l'alfabeto di
riferimento è quello delle cifre decimali,
\[
  \S = \set{0,1,\dots,9}, \qquad d = \abs{\S} = 10,
\]
ma la scelta è puramente espositiva: l'unica ipotesi che conta davvero è che l'alfabeto abbia
taglia costante, $d = \abs{\S} \in \Oh(1)$. In generale identifichiamo ogni carattere con la
cifra corrispondente in $\set{0,1,\dots,d-1}$, così che scrivere $P[i]$ o $T[j]$ dentro una
formula abbia senso.

Il problema è il consueto problema di pattern matching: determinare tutti gli \emph{shift
validi}, cioè gli $s \in \set{0,1,\dots,n-m}$ tali che $P = \sub{T}{s+1}{s+m}$.

\begin{figure}[H]\centering% kr-finestra-testo -- pag. 26
% Il testo T con la finestra corrente T[s+1..s+m] e la finestra successiva,
% traslata di una sola posizione; il pattern P ha la stessa lunghezza m.
\begin{tikzpicture}[
    >=Latex,
    et/.style={font=\footnotesize},
    graffa/.style={decorate,decoration={brace,amplitude=4pt,raise=1pt},
                   draw=black!70,line width=0.5pt}
  ]

  % ---- il testo T -----------------------------------------------------------
  % la finestra corrente, evidenziata sul testo
  \draw[line width=3.2pt,black!18] (1.00,0) -- (3.20,0);
  \draw[line width=0.9pt,black!80] (0,0) -- (8.40,0);
  \draw[line width=1.6pt,black]    (7.95,0) -- (8.40,0);
  \node[et,right=3pt] at (8.40,0) {$T$};

  \draw[black!60,line width=0.4pt] (1.00,-0.11) -- (1.00,0.11);
  \draw[black!60,line width=0.4pt] (3.20,-0.11) -- (3.20,0.11);
  \node[et,below=1pt] at (1.00,-0.11) {$s+1$};
  \node[et,below=1pt] at (3.20,-0.11) {$s+m$};

  % ---- finestra corrente  t_s = T[s+1..s+m] --------------------------------
  \draw[graffa] (1.00,0.10) -- (3.20,0.10);
  \node[et] at (2.10,0.60) {$\sub{T}{s+1}{s+m}$};

  % ---- finestra successiva: la stessa, traslata di una posizione ------------
  \draw[graffa] (1.45,1.40) -- (3.65,1.40);
  \node[et] at (2.55,1.90) {$\sub{T}{s+2}{s+m+1}$};

  % ---- lo scorrimento di una sola posizione --------------------------------
  \draw[->,black!70,line width=0.5pt] (3.20,1.05) -- (3.65,1.05);
  \node[et,right=3pt,black!70] at (3.65,1.05) {shift di una posizione};

  % ---- il pattern P, allineato sotto la finestra corrente ------------------
  \draw[black!35,dashed,line width=0.4pt] (1.00,-0.55) -- (1.00,-1.15);
  \draw[black!35,dashed,line width=0.4pt] (3.20,-0.55) -- (3.20,-1.15);
  \draw[line width=0.9pt,black!80] (1.00,-1.15) -- (3.20,-1.15);
  \draw[line width=1.6pt,black]    (2.85,-1.15) -- (3.20,-1.15);
  \node[et,right=3pt] at (3.20,-1.15) {$P$};

  % ---- la lunghezza comune e' m --------------------------------------------
  \draw[decorate,decoration={brace,mirror,amplitude=4pt,raise=1pt},
        draw=black!70,line width=0.5pt] (1.00,-1.25) -- (3.20,-1.25);
  \node[et] at (2.10,-1.72) {$m$};

\end{tikzpicture}

\caption{Il testo $T$, la finestra corrente $\sub{T}{s+1}{s+m}$ e la finestra successiva
(traslata di una posizione); il pattern $P$ ha la stessa lunghezza $m$ della finestra.}
\label{fig:kr-finestra-testo}
\end{figure}

\begin{osservazione}[Idea di partenza]
$\sub{P}{1}{m}$ può essere visto come un \emph{numero} scritto in base $d = \abs{\S}$: ogni
carattere è una cifra e la stringa non è altro che la sua rappresentazione posizionale. La
stessa lettura si applica, evidentemente, a ogni finestra di $T$ lunga $m$.
\end{osservazione}

Poniamo dunque:
\begin{itemize}
  \item $p$ è il valore di $\sub{P}{1}{m}$ letto in base $d$, cioè
    \[
      p \;=\; \sum_{i=1}^{m} P[i]\, d^{\,m-i}
        \;=\; P[1]d^{\,m-1} + P[2]d^{\,m-2} + \dots + P[m-1]d + P[m];
    \]
  \item $t_s$ è il valore di $\sub{T}{s+1}{s+m}$ letto in base $d$, cioè
    \[
      t_s \;=\; \sum_{i=1}^{m} T[s+i]\, d^{\,m-i}.
    \]
\end{itemize}

Poiché la scrittura posizionale in base $d$ di stringhe di lunghezza \emph{fissa} $m$ è
iniettiva, vale
\[
  s \text{ è uno shift valido} \iff p = t_s,
\]
e il confronto fra stringhe si è trasformato in un confronto fra interi.

\subsection{Lo schema di partenza}

Lo schema di algoritmo che ne discende è il seguente.
\begin{enumerate}
  \item Inizializzo $p$ e $t_0$.
  \item Calcolo $t_{s+1}$ a partire da $t_s$ con uno \emph{shift} e una \emph{somma}: dalla
    finestra esce il carattere $T[s+1]$ ed entra $T[s+m+1]$, quindi basta togliere la cifra
    più significativa, moltiplicare per $d$ (lo shift) e sommare la nuova cifra meno
    significativa.
  \item A ogni passo confronto $p$ con $t_s$: se sono uguali segnalo un'occorrenza,
    altrimenti eseguo lo shift e passo a $s+1$.
\end{enumerate}

\begin{figure}[H]\centering% kr-schema-loop -- pag. 26
% Lo schema "ingenuo" di Karp--Rabin: un unico test p = t_s, ripetuto per tutti
% gli shift, con aggiornamento incrementale della finestra; l'annotazione in basso
% e' il punto debole (il test non costa O(1)).
\begin{tikzpicture}[
    >=Latex,
    et/.style={font=\footnotesize},
    filo/.style={draw=black!75,line width=0.5pt},
    blocco/.style={draw=black!75,line width=0.7pt,minimum width=3.6cm,
                   minimum height=1.1cm,inner sep=4pt,font=\small}
  ]

  % ---- il blocco di test ----------------------------------------------------
  \node[blocco] (test) at (0,0) {test\quad $p = t_s$};

  % ---- uscite a ventaglio dal lato destro -----------------------------------
  \coordinate (fork) at (1.80,0);
  \node[et,anchor=west] (occ)   at (3.35, 0.80) {occorrenza};
  \node[et,anchor=west] (shift) at (3.35,-0.80) {shift e somma};

  \draw[filo,->] (fork) -- (3.25, 0.80);
  \draw[filo,->] (fork) -- (3.25,-0.80);
  \node[et] at (2.30, 0.62) {s\`i};
  \node[et] at (2.30,-0.62) {no};

  % ---- segnalata l'occorrenza, la scansione prosegue comunque ---------------
  \draw[filo,->,black!45,densely dashed] (occ.south) -- (shift.north);

  % ---- retroazione: il ciclo su tutti gli shift ------------------------------
  \draw[filo,->]
      (shift.south) to[out=-90,in=0] (1.10,-2.95)
                    to[out=180,in=-90] (-3.45,-0.55)
                    to[out=90,in=180] (test.west);
  \node[et,below=2pt,black!70] at (1.10,-2.95) {ciclo su $s = 0,1,\dots,n-m$};

  % ---- annotazione: il test non e' a costo unitario -------------------------
  \draw[filo,black!60] (-0.50,-0.55) -- (-1.30,-1.05);
  \node[anchor=north,align=center,font=\small] at (0.75,-1.05)
    {{\large\bfseries NON} prende tempo $\Oh(1)$!\\[2pt]
     \footnotesize Non si pu\`o usare il criterio di costo unitario};

\end{tikzpicture}

\caption{Lo schema di Karp--Rabin: un unico test $p = t_s$ ripetuto per tutti gli shift, con
aggiornamento incrementale della finestra.}
\label{fig:kr-schema-loop}
\end{figure}

\begin{osservazione}[Il punto debole]
Il test $p = t_s$ \textbf{non} prende tempo $\Oh(1)$: su questi dati non si può usare il
criterio di costo unitario. Infatti $p$ e i $t_s$ sono numeri di $m$ cifre in base $d$, cioè
richiedono $\Th(m \log d) = \Th(m)$ bit; non appena $m$ supera l'ordine di $\log n$ non stanno
più in una parola di macchina, e allora ogni confronto (come ogni aggiornamento) costa
$\Th(m)$ anziché $\Oh(1)$. Ripetendolo per tutti gli $n-m+1$ shift si ritorna a
$\Th\bigl((n-m+1)\,m\bigr)$: nessun guadagno rispetto all'algoritmo naive.
\end{osservazione}

\section{Il rimedio: lavorare modulo \texorpdfstring{$q$}{q}}

\scan{27}
Il problema, insomma, è che \emph{$p$ e i $t_s$ sono troppo grandi}. La via d'uscita è non
manipolarli mai, ma manipolare soltanto i loro resti
\[
  p \mod q, \qquad t_s \mod q,
\]
dove $q$ è un parametro scelto in modo da garantire che $p \mod q$ (e analogamente ogni
$t_s \mod q$) sia manipolabile in tempo $\Oh(1)$.

\begin{proposizione}[Il test modulare è un filtro a un lato]\label{prop:kr-filtro}
Per ogni intero $q \ge 1$ e per ogni shift $s$ vale l'implicazione
\[
  p \mod q \neq t_s \mod q \implies p \neq t_s,
\]
e dunque $s$ non è uno shift valido. L'implicazione inversa, invece, \emph{non} vale in
generale: non appena $q < d^{\,m}$ esistono un pattern e una finestra di testo con
$p \neq t_s$ e $p \mod q = t_s \mod q$, cioè
\[
  p \mod q = t_s \mod q \;\nRightarrow\; p = t_s .
\]
\end{proposizione}

\begin{dimostrazione}
La prima implicazione è la contronominale del fatto che interi uguali hanno lo stesso resto:
se $p = t_s$ allora $p \mod q = t_s \mod q$. Dunque, se i due resti differiscono,
necessariamente $p \neq t_s$ e quindi --- essendo la scrittura in base $d$ di una stringa di
lunghezza fissa $m$ iniettiva --- lo shift $s$ non è valido.

La seconda non vale: $p$ e $t_s$ variano in $\set{0,1,\dots,d^{\,m}-1}$, che ha $d^{\,m}$
elementi, mentre i resti modulo $q$ sono soltanto $q$; se $q < d^{\,m}$, per il principio dei
cassetti due valori distinti finiscono sullo stesso resto. Basta dunque considerare il caso in
cui $p \neq t_s$ ma
\[
  \abs{p - t_s} \mod q = 0,
\]
cioè $q \mid p - t_s$ con $p \neq t_s$: allora $p$ e $t_s$ hanno lo stesso resto modulo $q$
pur essendo diversi, e il test modulare non li distingue. (Se $q \ge d^{\,m}$ il test è
esatto, ma è proprio il caso che si vuole evitare: i resti sarebbero grandi quanto $p$ e
$t_s$, e si tornerebbe al problema di partenza.)
\end{dimostrazione}

\begin{osservazione}
Il confronto fra resti è quindi un filtro a un solo lato: quando i resti differiscono lo
shift si può scartare con certezza (mai un falso negativo), ma quando coincidono non si
conclude nulla, perché $p \mod q = t_s \mod q$ è solo condizione \emph{necessaria} perché $s$
sia uno shift valido. I falsi positivi sono esattamente le posizioni con $p \neq t_s$ e
$q \mid p - t_s$; ognuna di esse va smascherata con una verifica esplicita.
\end{osservazione}

\paragraph{Problemi.}
\begin{enumerate}
  \item[i)] come scelgo $q$?
  \item[ii)] come calcolo e aggiorno $t_s \mod q$?
  \item[iii)] cosa faccio quando $p \mod q = t_s \mod q$?
\end{enumerate}

\begin{dubbio}[Annotazione a matita, in parte cancellata]
A metà pagina compare, a matita e quasi del tutto cancellata, la fattorizzazione
\[
  q = d^{\,n}\cdot a
  \qquad\longleftrightarrow\qquad
  \begin{cases}
    q' \mod d^{\,n} & \dfrac{1}{d^{\,n}}\\[2ex]
    q'' \mod a      & \dfrac{1}{a}
  \end{cases}
  \qquad\qquad \frac{1}{d^{\,n}}\cdot\frac{1}{a},
\]
accanto a uno schizzo (anch'esso a matita e cancellato) di un segmento con alcune tacche e un
cappio, privo di etichette. L'esponente, letto $n$, non ha nulla a che vedere con
$n = \abs{T}$, e la lettera $a$ è qui un cofattore, non il carattere entrante dello
pseudocodice.

La lettura più plausibile è che si tratti di una verifica della probabilità $1/q$ di
collisione: se $q$ si fattorizza come $d^{\,n}\cdot a$ con $\gcd(d^{\,n},a)=1$, per il
teorema cinese del resto la condizione $p \equiv t_s \pmod q$ equivale alla congiunzione di
$p \equiv t_s \pmod{d^{\,n}}$ e $p \equiv t_s \pmod a$, eventi ``indipendenti'' di
probabilità $1/d^{\,n}$ e $1/a$, il cui prodotto è appunto $1/q$. Si noti però che un $q$ di
quella forma sarebbe una pessima scelta: la riduzione modulo $d^{\,n}$ (con $d = \abs{\S}$ la
base) legge soltanto gli ultimi $n$ caratteri della finestra. È anche per evitare queste
degenerazioni che si prenderà $q$ primo.
\end{dubbio}

\section{L'algoritmo}

Rimandata la scelta di $q$, l'algoritmo prende la forma seguente. Da qui in avanti le
variabili che si chiamano $p$ e $t_s$ contengono i \emph{resti} $p \mod q$ e $t_s \mod q$,
non più i due interi grandi: è l'abuso di notazione consueto, e non crea ambiguità perché
$p$ e $t_s$ nel loro significato originario non vengono mai calcolati.

\begin{algorithm}[H]
\caption{\textsc{Karp--Rabin}$(T,P,d,q)$}
\label{alg:karp-rabin}
\begin{algorithmic}[1]
  \State $m \gets P.\mathit{length}$
  \State $n \gets T.\mathit{length}$
  \State $h \gets d^{\,m-1} \mod q$
  \State $p \gets 0$
  \State $t_0 \gets 0$
  \For{$i \gets 1$ \textbf{to} $m$} \Comment{regola di Horner, modulo $q$}
    \State $p   \gets (d\cdot p   + P[i]) \mod q$
    \State $t_0 \gets (d\cdot t_0 + T[i]) \mod q$
  \EndFor
  \For{$s \gets 0$ \textbf{to} $n-m$}
    \If{$p = t_s$} \Comment{filtro}
      \If{$\sub{P}{1}{m} = \sub{T}{s+1}{s+m}$} \Comment{verifica}
        \State \textbf{match} allo shift $s$
      \EndIf
    \EndIf
    \If{$s < n-m$} \Comment{aggiornamento della finestra}
      \State $r \gets (T[s+1]\cdot h) \mod q$
      \State $a \gets T[s+m+1]$
      \State $t_{s+1} \gets \bigl(d\,(t_s - r) + a\bigr) \mod q$
    \EndIf
  \EndFor
\end{algorithmic}
\end{algorithm}

Negli appunti una graffa raccoglie le righe $1$--$9$ e le etichetta $\Oh(m)$: è tutto il
pre-processing, cioè le due valutazioni con la regola di Horner, una moltiplicazione e una
somma per carattere, eseguite direttamente modulo $q$. Anche la riga $3$ rientra nel budget:
$h = d^{\,m-1} \mod q$ si calcola con l'esponenziazione veloce in $\Oh(\log m)$ operazioni, o
banalmente in $\Oh(m)$.

Le tre righe interne al secondo ciclo sono l'aggiornamento incrementale della finestra
descritto al passo $2$ dello schema: $r$ è il contributo della cifra che esce, $a$ è la cifra
che entra, e in forma chiusa
\[
  t_{s+1} = \bigl(d\,(t_s - T[s+1]\cdot h) + T[s+m+1]\bigr) \mod q,
  \qquad h = d^{\,m-1} \mod q,
\]
cioè $\Oh(1)$ operazioni per shift; la guardia $s < n-m$ serve solo a non leggere $T[s+m+1]$
oltre la fine del testo. Che tutto questo si possa fare direttamente sui resti dipende dal
fatto che la riduzione modulo $q$ commuta con somma, differenza e prodotto.

Il confronto $p = t_s$ è il filtro della Proposizione~\ref{prop:kr-filtro}, e il confronto
carattere per carattere $\sub{P}{1}{m} = \sub{T}{s+1}{s+m}$ è la verifica che esso rende
necessaria --- la risposta al problema iii): al più $m$ confronti, ed è il punto in cui si
concentrerà l'analisi di complessità.

\section{La scelta di \texorpdfstring{$q$}{q}}

\scan{28}
Resta aperto il problema i). Il parametro $q$ va preso \emph{primo}, e più grande è, meglio
è: il che è ovvio, dato che la probabilità di un falso positivo sarà $1/q$. Dall'altro lato
$q$ non può crescere quanto si vuole, perché tutto l'impianto si regge sull'ipotesi che le
operazioni modulo $q$ costino $\Oh(1)$: $q$ (e con esso $d\cdot q$) deve stare in una parola
di macchina, cioè occupare $\Oh(\log n)$ bit. I due vincoli sono compatibili --- l'analisi
della prossima sezione mostrerà che basta $q$ dell'ordine di $m$, e $m \le n$ --- e resta da
capire se un primo della taglia voluta esista davvero e se sia facile trovarlo: per questo
useremo il teorema sulla distribuzione dei primi.

Il modulo mappa gli interi in $\set{0,\dots,q-1}$. In modo uniforme? Sì --- o, più
onestamente, \emph{lo assumiamo}: è esattamente questa ipotesi che permette di dire che per
uno shift non valido la congruenza $p \equiv t_s \pmod{q}$ si presenta con probabilità $1/q$.

\section{Analisi di complessità}

Prima dell'analisi vera e propria, negli appunti si apre a margine una parentesi
terminologica.

\begin{osservazione}[Smoothed complexity]
La \emph{smoothed complexity} ``liscia'' la superficie di complessità. Ci sono algoritmi, come
QuickSort, che benché su carta siano $\Oh(n^{2})$, nella pratica sono solitamente più veloci:
l'analisi smoothed serve proprio a rendere conto di questo scarto, facendo pesare di meno i
casi peggiori isolati --- gli \emph{outlier}.
\end{osservazione}

\begin{dubbio}[Ultima riga della parentesi]
L'ultima riga della parentesi si legge ``nella smoothed complexity si \dots{} gruppi maggiori
ridotti dai outlier'', con due verbi separati da una barra che non è stato possibile
decifrare. La lettura qui adottata --- si valuta il costo su \emph{gruppi} (intorni) di input
anziché sul singolo input, così che il contributo degli outlier risulti ridotto --- è
congetturale.
\end{dubbio}

Veniamo all'analisi. Siano
\begin{itemize}
  \item $v$ il numero di posizioni in cui $p \equiv t_s \pmod{q}$ \emph{e} $P$ occorre
        davvero (occorrenze \emph{valide});
  \item $k$ il numero di posizioni in cui $p \equiv t_s \pmod{q}$ \emph{ma} $P$ non occorre
        (occorrenze \emph{spurie}).
\end{itemize}
In entrambi i casi scatta il confronto carattere per carattere, che costa $\Oh(m)$.

Da qui due conseguenze immediate:
\begin{itemize}
  \item se $v$ è alto, Karp--Rabin costa $\Oh(m\cdot n)$, esattamente quanto l'algoritmo
        naive;
  \item assumo dunque che $v$ sia $\Oh(1)$.
\end{itemize}

\begin{osservazione}[Annotazione a matita]
Accanto all'ipotesi $v = \Oh(1)$ il quaderno annota, a matita: ``non è vero per testi
completamente random''. L'osservazione va però capovolta: se $T$ è davvero uniforme e casuale
su un alfabeto di taglia $d$, il numero atteso di occorrenze di $P$ è $(n-m+1)/d^{\,m}$, che
è praticamente nullo non appena $d^{\,m} \gg n$, cioè per $m \gtrsim \log_d n$: in quel regime
l'ipotesi $v = \Oh(1)$ è ampiamente soddisfatta. (Per $m$ molto piccolo il valore atteso di
$v$ può essere $\Th(n)$ anche su testo casuale, ma allora $m\,\E[v]$ resta $\Oh(n)$ e la
complessità non ne risente.) È sui testi \emph{ripetitivi} che l'ipotesi cade: con $P = a^{m}$
e $T = a^{n}$ si ha $v = n-m+1$, e in quel caso nessun algoritmo che verifichi tutte le
occorrenze può scendere sotto $\Th\bigl((n-m+1)\,m\bigr)$.
\end{osservazione}

Quanto a $k$, sono costretto a spendere $\Oh(m)$ ogni volta che
\[
  p - t_s \equiv 0 \pmod{q},
\]
cosa che --- sotto l'ipotesi di uniformità di poco fa --- avviene con probabilità $1/q$;
dunque
\[
  \E[k] \;\le\; \frac{n-m+1}{q} \;=\; \Oh\!\left(\frac{n}{q}\right).
\]

Mettendo tutto insieme (il primo addendo è la scansione del testo con l'aggiornamento
incrementale della finestra, il secondo sono le verifiche esplicite) si ottiene, in valore
atteso,
\begin{align*}
  \Oh(n) + \Oh\bigl(m\,(v+k)\bigr)
    &= \Oh(n) + \underbrace{\Oh(m\,v)}_{v \,=\, \Oh(1)} + \Oh(m\,k) \\[2pt]
    &= \Oh(n) + \Oh(m) + \Oh\!\left(\frac{m\cdot n}{q}\right),
\end{align*}
e con $q \gtrsim m$ il tutto è \emph{lineare}: $\Oh(n+m) = \Oh(n)$.

\begin{osservazione}[Che cosa è davvero aleatorio]
Il conto appena fatto è un conto \emph{in valore atteso}, e l'aspettazione è presa rispetto
all'ipotesi di uniformità dei resti. Con $q$ fissato una volta per tutte quell'ipotesi è una
comodità di analisi e non un teorema: un testo costruito ad arte può far collidere tutti gli
shift e riportare il costo a $\Th\bigl((n-m+1)\,m\bigr)$. La versione rigorosa sceglie $q$ primo \emph{a caso}
in un intervallo opportuno, ed è allora la scelta di $q$ --- non il testo --- a fornire
l'aleatorietà su cui si prende il valore atteso.
\end{osservazione}
